\chapter{Introducción y Objetivos}
\label{chap:introduccion}

El diseño de etapas analógicas discretas con transistores bipolares de juntura (BJT) requiere seleccionar la
topología circuital adecuada según los requerimientos específicos de adaptación de impedancias, ganancia y
ancho de banda. En aplicaciones electroacústicas de instrumentación y acondicionamiento de señales, es habitual
enfrentar transductores de alta impedancia interna acoplados a cargas de valor moderado mediante líneas de
transmisión extensas, situación donde la correcta elección de la configuración del transistor resulta crítica.

\section{Objetivos del Trabajo Práctico}

El presente trabajo de laboratorio persigue los siguientes objetivos técnicos y formativos:
\begin{enumerate}
  \item Analizar un requerimiento práctico de acondicionamiento de señal y seleccionar justificadamente la
    topología con transistor BJT que mejor satisfaga las especificaciones del sistema.
  \item Diseñar, implementar y verificar experimentalmente una etapa amplificadora para cumplir con las
    especificaciones de ganancia de tensión ($A_v$), ganancia de corriente ($A_i$), impedancia de entrada ($Z_i$),
    impedancia de salida ($Z_o$) y excursión dinámica libre de distorsión.
  \item Emplear herramientas de simulación computarizada en LTspice y métodos analíticos para predecir el
    comportamiento estático en continua y dinámico en alterna.
  \item Normalizar los componentes calculados a valores comerciales de la serie estándar E12/E24 y verificar el
    impacto sobre el punto estático de reposo.
  \item Trazar y analizar las rectas de carga estática de continua y dinámica de alterna, evaluando los márgenes
    reales de excursión lineal.
  \item Montar físicamente el circuito en placa de circuito impreso (PCB) en laboratorio y relevar mediante
    instrumental de banco las magnitudes estáticas y dinámicas.
  \item Contrastar de forma cuantitativa los resultados analíticos, las simulaciones circuitales y las mediciones
    de laboratorio, evaluando fuentes de incertidumbre y errores experimentales.
\end{enumerate}

\section{Resultados de Aprendizaje Vinculados}

La realización de esta práctica consolida los siguientes resultados de aprendizaje de la asignatura Electrónica
Aplicada I:
\begin{itemize}
  \item RA1: Utilizar instrumentos de laboratorio de generación, visualización y medición de señales (osciloscopio,
    generador de funciones, multímetros) para analizar el funcionamiento de circuitos electrónicos analógicos.
  \item RA3: Desarrollar amplificadores discretos clase A con BJT utilizando conceptos de estabilidad térmica para
    aplicaciones en audiofrecuencia.
\end{itemize}

\section{Materiales e Instrumental}

Para el desarrollo del trabajo se emplearon los siguientes elementos de hardware y laboratorio:
\begin{itemize}
  \item Transistor bipolar NPN de silicio de uso general tipo BC547B / BC548.
  \item Resistores de película de carbón de $0.25\W$ según dimensionamiento ($47\kohm$, $82\kohm$, $1\kohm$, $10\kohm$,
    y resistores de prueba para caracterización de impedancias).
  \item Capacitores electrolíticos y de película de poliéster para acoplo de señal ($22\uF$ / $47\uF$).
  \item Placa de circuito impreso (PCB) de prueba para montaje discreto.
  \item Generador de funciones de laboratorio y osciloscopio de doble canal con puntas atenuadoras calibradas.
  \item Multímetros digitales de banco y portátiles.
  \item Fuente de alimentación de corriente continua regulada de $0\V$ a $30\V$.
  \item Computadora personal con entorno de simulación LTspice XVII.
\end{itemize}
